% ============================================
% Chapter 8: System Integration
% ============================================
% NEW chapter --- Adil writes this

\chapter{System Integration}

\section{Introduction}

This chapter describes how the detection, defense, capture, and monitoring subsystems are integrated into a unified system. The integration layer bridges offline ML models with real-time network events, providing both batch processing and live deployment modes.

\section{Integration Architecture}

\subsection{Bridge Module}

The Bridge module (\texttt{bridge/bridge.py}) serves as the central integration point, connecting Nmap scan results and live captures to the ML detection pipeline and the Flask dashboard.

\begin{figure}[H]
\centering
\begin{tikzpicture}[
    node distance=1.5cm,
    block/.style={rectangle, draw=cyan, fill=darkbg!5, text width=3cm,
                  minimum height=0.9cm, align=center, rounded corners=4pt,
                  font=\small\sffamily, line width=1pt},
    arrow/.style={-{Stealth[length=3mm]}, thick, cyan!70!black},
    lbl/.style={font=\tiny\sffamily, color=codegray, midway}
]
    \node[block] (nmap) {Nmap XML\\Parser};
    \node[block, right=2cm of nmap] (bridge) {Bridge\\Module};
    \node[block, above right=1cm and 1.5cm of bridge] (ml) {ML\\Pipeline};
    \node[block, below right=1cm and 1.5cm of bridge] (dash) {Flask\\Dashboard};
    \node[block, below=1.5cm of bridge] (local) {Local\\JSONL Log};

    \draw[arrow] (nmap) -- node[lbl,above]{features} (bridge);
    \draw[arrow] (bridge) -- node[lbl,above]{predictions} (ml);
    \draw[arrow] (bridge) -- node[lbl,right]{POST /api/alert} (dash);
    \draw[arrow] (bridge) -- node[lbl,left]{fallback} (local);
\end{tikzpicture}
\caption{Bridge module data flow: Nmap XML $\rightarrow$ ML prediction $\rightarrow$ dashboard + local log.}
\label{fig:bridge_flow}
\end{figure}

\subsection{Nmap Parser}

The Nmap Parser (\texttt{bridge/nmap\_parser.py}) converts Nmap XML output files into feature vectors compatible with the ML pipeline. It extracts relevant fields from \texttt{<port>} and \texttt{<service>} elements and maps them to the 11-feature schema.

\subsection{Dual Logging}

All subsystems implement dual logging to ensure data availability regardless of dashboard connectivity:

\begin{itemize}
    \item \textbf{Local JSONL}: Events appended to \texttt{detection\_logs.json} and \texttt{deception\_logs.json} in the \texttt{detection/data/} directory
    \item \textbf{HTTP POST}: Events pushed to the Flask dashboard via \texttt{/api/alert}, \texttt{/api/block}, and \texttt{/api/honeypot\_event} endpoints
\end{itemize}

If the dashboard is unreachable, events are written locally as a fallback (no data loss).

\section{Dashboard}

\subsection{Backend}

The Flask dashboard (\texttt{dashboard/app.py}) provides REST API endpoints and Socket.IO real-time event streaming. It maintains JSONL log files for detection and deception events.

\subsection{Frontend}

The SOC dashboard frontend (\texttt{dashboard/templates/index.html}) displays:
\begin{itemize}
    \item Active alerts with confidence scores and severity levels
    \item Blocked IP list with unblock controls
    \item Honeypot interaction events
    \item Real-time event stream via Socket.IO
\end{itemize}

\section{Deployment Modes}

The system supports three deployment modes via the unified entry point (\texttt{run\_demo.py}):

\begin{table}[H]
\centering
\caption{Deployment modes}
\label{tab:deploy_modes}
\begin{tabular}{llp{6cm}}
\toprule
\textbf{Mode} & \textbf{Flag} & \textbf{Description} \\
\midrule
Offline & \texttt{--offline} & Process Nmap XML files through ML pipeline; no live traffic capture \\
Live & \texttt{--live} & Real-time Scapy capture + ML classification + defense response \\
Dashboard only & \texttt{--dashboard-only} & Start Flask dashboard without ML or capture modules \\
\bottomrule
\end{tabular}
\end{table}

\section{Conclusion}

The integration layer unifies all subsystems into a coherent pipeline. The Bridge module handles offline and batch processing, while the live mode provides real-time detection and response. Dual logging ensures no data loss regardless of dashboard availability.
